As an example of the conjecture, consider the following quiver.
\[
\xymatrixrowsep{15pt}\xymatrixcolsep{10pt}
\xymatrix{%begin xy matrix
Q:&& 2\ar[dl]_\alpha\\
&1\ar[rr]^\gamma &&3\ar[lu]_\beta
	}%end xy matrix
\]
with potential $W=\alpha\beta\gamma$. Potentials are linear combinations of oriented cycles in the quiver and are only well-defined up to cyclic permutation. The \emph{Jacobian algebra} $J(Q,W)$ is defined to be the path algebra $KQ$ of $Q$ over a fixed field $K$ modulo the partial derivatives of $W$ with respect to each arrow~\cite{DWZ1}. In this case these relations~are:
\[
	\partial_\alpha W:=\beta\gamma=0,\quad \partial_\beta W:=\gamma\alpha=0,\quad \partial_\gamma W:=\alpha\beta=0.
\]
In other words, $\Lambda=J(Q,W)=KQ/\rad ^2KQ$. 

As we explain in the Appendix, there are three ``tilted algebras'' associated to $\Lambda$. They are given by ``cutting'' the quiver, \ie, by deleting one of the three edges of $Q$. Any two of the quivers obtained by cutting are isomorphic. So, we delete edge $\gamma$. We are left with the following quiver:
\[
	Q^\delta:\quad 1\xleftarrow\alpha 2\xleftarrow\beta 3
\]
with one relation $\alpha\beta=0$ since this is the only relation of $Q$ supported on the subquiver $Q^\delta$. The resulting algebra $C=KQ^\delta/(\alpha\beta)$ is a tilted algebra of type $A_3$ with 5 indecomposable modules with dimension vectors $(0,0,1), (0,1,1), (0,1,0), (1,1,0), (1,0,0)$. Conjecture~\ref{conjectural description of maximal MGSs} then states that a longest maximal green sequence for $Q$ has these vectors as the $c$-vectors of the mutations as we now explain and thus has length 5.

Maximal green sequences (MGS) are defined only in terms of the quiver $Q$ as follows. The \emph{exchange matrix} $B$ of $Q$ is defined to be the skew-symmetric matrix with $ij$th entry $b_{ij}$ equal to the number of arrows $i\to j$ in the quiver $Q$ minus the number of arrows $j\to i$. %For this example,
\[
	B=\mat{0 & -1 & 1\\
	1 & 0 & -1\\
	-1 & 1 & 0}
\]
The \emph{extended exchange matrix} $\tilde B$ is given by putting the $3\times 3$ identity matrix $I_3$ below~$B$. By Theorem~\ref{thm: MSG for An} there is a MGS of length 5 which mutates at the \hbox{``$c$-vectors''} which are the dimension vectors of the $C$-modules as follows.
\begin{gather*}
	\tilde B=\mat{0 & -1 & 1\\
	1 & 0 & -1\\
	-1 & 1 & 0\\
	\hline % first matris
	1 & 0 & \color{green}0\\
	0 & 1 & \color{green}0\\
	0 & 0 & \color{green}1} 
	\xrightarrow{\mu_3}
	 \mat{0 & 0 & -1\\
	0 & 0 & 1\\
	1 & -1 & 0\\
	\hline % second matrix
	1 & \color{green}0 & 0\\
	0 & \color{green}1 & 0\\
	0 & \color{green}1 & -1} 
	\xrightarrow{\mu_2} 
	 \mat{0 & 0 & -1\\
	0 & 0 & -1\\
	1 & 1 & 0\\
	\hline % third matrix
	1 & 0 & \color{green}0\\
	0 & -1 & \color{green}1\\
	0 & -1 & \color{green}0} 
	\xrightarrow{\mu_3} 
	\mat{0 & 0 & 1\\
	0 & 0 & 1\\
	-1 & -1 & 0\\
	\hline % fourth matrix
	\color{green}1 & 0 & 0\\
	\color{green}1 & 0 & -1\\
	\color{green}0 & -1 & 0} 
\\
%	\]
%	\[
	\xrightarrow{\mu_1} 
	\mat{0 & 0 & -1\\
	0 & 0 & 1\\
	1 & -1 & 0\\
	\hline % fifth matrix
	-1 & 0 & \color{green}1\\
	-1 & 0 & \color{green}0\\
	0 & -1 & \color{green}0}	
	\xrightarrow{\mu_3} 
	\mat{0 & -1 & 1\\
	1 & 0 & -1\\
	-1 & 1 & 0\\
	\hline % sixth matrix
	0 & 0 & -1\\
	-1 & 0& 0\\
	0 & -1 & 0}
\end{gather*}
The \emph{$c$-vectors} are the columns of the bottom half of these exchange matrices. The Fomin--Zelevinsky mutation $\mu_k$ is given by applying the following general rule. We add $b_{ik}|b_{kj}|$ to $b_{ij}$ whenever $b_{ik}$ and $b_{kj}$ have the same sign, then we change the signs of all $b_{ik},b_{kj}$~\cite{FZ4}. The mutation $\mu_k$ is \emph{green} if the $k$th $c$-vector is positive, i.e., all its entries are nonnegative. The last matrix has no positive $c$-vectors. So, the green mutation sequence must stop. This is therefore a \emph{maximal green sequence}. The standard notation for this mutation sequence is $(3,2,3,1,3)$. However, we prefer to label the sequence with its $c$-vectors because of their representation theoretic interpretation.


The $c$-vectors of this MGS are the dimension vectors of the $C$-modules in the order $M_1,\ldots,M_5$ so that $\Hom_{C}(M_i,M_j)=\Hom_{\Lambda}(M_i,M_j)=0$ for $i<j$. The sixth $\Lambda$-module $M_6$ with $\undim M_6=(1,0,1)$ cannot be inserted into this sequence since $\Hom_\Lambda(M_1,M_6)\neq 0$ and $\Hom_\Lambda(M_6,M_5)\neq 0$. So, the sequence $M_1,\ldots,M_5$ is a \emph{complete forward hom-orthogonal sequence}.

